\section{Markov Networks}
\label{sec:markov-networks}

This section introduces the model class used as the final decision
layer of the predictor: Markov Networks (MNs), also known as Markov
Random Fields. An MN is a graphical model in which dependencies
between output variables are encoded by an undirected graph, and the
joint score of a configuration decomposes into local terms attached
to nodes and edges. Both structural variants considered in this
thesis --- linear chains and general graphs --- are special cases of
the formalism developed below. The setting is \emph{discriminative}:
the score $f(x, y)$ conditions on the input $x$ and the model does
not attempt to estimate $p(x)$, distinguishing it from generative
graphical models such as Hidden Markov Models in which the joint
distribution $p(x, y)$ is modeled and the input is treated as a
random variable. A comprehensive treatment of graphical models is
given by~\cite{Koller2009}.

\subsection{Graph, variables, and factorization}

Let $G = (V, E)$ be an undirected graph with $N = |V|$ nodes. Each
node $v \in V$ carries a discrete random variable $y_v$ taking values
in a finite label set $\mathcal{Y} = \{1, \ldots, K\}$. A joint
configuration is therefore an element $y = (y_v)_{v \in V} \in
\mathcal{Y}^{N}$.

For the tasks considered in this thesis, the score of a configuration
decomposes over the nodes and edges of $G$:
\begin{equation}
    f(x, y) \;=\; \sum_{v \in V} f_v(x, y_v)
    \;+\; \sum_{\{v, v'\} \in E} f_{vv'}(y_v, y_{v'}),
    \label{eq:mn-score}
\end{equation}
where $f_v$ is a \emph{unary} term coupling the observation $x$ to the
label of node $v$, and $f_{vv'}$ is a \emph{pairwise} term encoding
the compatibility of the labels at neighbouring nodes. Equation
\eqref{eq:mn-score} describes a \emph{pairwise} MN; higher-order
cliques are possible in general, but the pairwise form is sufficient
for both the linear-chain and the Sudoku models used here.

The undirected graph $G$ encodes a conditional-independence
structure: given the labels of its neighbours, $y_v$ is conditionally
independent of all other variables. This is the defining
\emph{Markov property} from which the model class takes its name.

\subsection{Linear parameterization}
\label{subsec:linear-param}

The scoring functions in Eq.~\eqref{eq:mn-score} are assumed to be
linear in a shared parameter vector $w \in \mathbb{R}^d$. Given fixed
feature maps
$\psi_v\colon \mathcal{X} \times \mathcal{Y} \to \mathbb{R}^d$ and
$\psi_{vv'}\colon \mathcal{Y} \times \mathcal{Y} \to \mathbb{R}^d$, the
unary and pairwise terms take the form
\begin{equation}
    f_v(x, y_v) = \langle w, \psi_v(x, y_v) \rangle,
    \qquad
    f_{vv'}(y_v, y_{v'}) = \langle w, \psi_{vv'}(y_v, y_{v'}) \rangle.
\end{equation}
Defining the \emph{joint feature map}
\begin{equation}
    \psi(x, y) \;=\; \sum_{v \in V} \psi_v(x, y_v)
    \;+\; \sum_{\{v, v'\} \in E} \psi_{vv'}(y_v, y_{v'})
    \label{eq:joint-feature}
\end{equation}
then yields the compact expression
\begin{equation}
    f(x, y) \;=\; \langle w, \psi(x, y) \rangle.
    \label{eq:linear-score}
\end{equation}
Equation~\eqref{eq:linear-score} is the starting point for the
max-margin learning theory of Section~\ref{sec:learning-mn}. The
joint-feature-map parameterization is standard in structured output
learning~\cite{Taskar2004, Tsochantaridis2005}. The architectures
used in this thesis depart from strict linearity
(Section~\ref{sec:nn-features}): the unary potential is produced by
a neural network with parameters $\theta$, and the pairwise
potential is a learned matrix $W \in \mathbb{R}^{K \times K}$ shared
across edges. The abstract notation $\langle w, \psi(x, y)\rangle$
is retained for the derivations of Section~\ref{sec:learning-mn};
from Section~\ref{subsec:training-procedure} onwards the parameters
are referenced as $(\theta, W)$ explicitly.

\subsection{Special cases}

Two instances of the general model above are used throughout the
thesis.

\paragraph{Linear-chain MN.}
For sequence labeling of length $T$, the graph is the linear chain
\begin{equation}
    V = \{0, 1, \ldots, T - 1\},
    \qquad
    E = \{\{t-1, t\} : t = 1, \ldots, T - 1\}.
\end{equation}
The score in Eq.~\eqref{eq:mn-score} reduces to
\begin{equation}
    f(x, y) \;=\; \sum_{t=0}^{T-1} f_t(x, y_t)
    \;+\; \sum_{t=1}^{T-1} f_{t-1, t}(y_{t-1}, y_t),
    \label{eq:chain-score}
\end{equation}
which is the discriminative, score-based form of a linear-chain
conditional random field~\cite{Lafferty2001}. MAP inference on chains
admits an exact $O(T \cdot K^2)$ dynamic-programming solution,
described in Section~\ref{sec:map-inference}.

\paragraph{Sudoku MN.}
For a $9 \times 9$ Sudoku puzzle, $V$ contains the 81 cells and
$\mathcal{Y} = \{1, \ldots, 9\}$. The edge set $E$ connects every pair
of cells that share a row, a column, or a $3 \times 3$ box ---
i.e.~every pair of cells that must take distinct values. Each cell
has $20$ such neighbours ($8$ in its row, $8$ in its column, and $4$
further cells in its $3 \times 3$ box), giving $|V| = 81$ and
$|E| = 810$. The pairwise potentials $f_{vv'}$ encode the \emph{all-different}
constraint along these edges. In principle they can be set as a hard
constraint ($f_{vv'}(a, a) = -\infty$ for all $a \in \mathcal{Y}$);
the architectures used in this thesis instead learn the pairwise
matrix $W \in \mathbb{R}^{K \times K}$ from data with no hard
constraint imposed (Section~\ref{sec:nn-features}).
The unary potentials $f_v(x, y_v)$ carry the evidence for cell $v$
coming from the input (a digit for symbolic input, an image for
visual input).

The Sudoku graph contains many cycles, so exact MAP inference is
NP-hard in general. Section~\ref{sec:map-inference} discusses
approximate MAP inference on such cyclic graphs.